\documentclass[11pt,a4paper]{article}
\usepackage[utf8]{inputenc}
\usepackage[T1]{fontenc}
\usepackage[indonesian]{babel}
\usepackage[margin=2.2cm]{geometry}
\usepackage{xcolor,listings,enumitem,booktabs,hyperref,parskip}
\definecolor{dawetblue}{HTML}{16324F}
\definecolor{codebg}{HTML}{F1F4F7}
\lstset{basicstyle=\ttfamily\small,backgroundcolor=\color{codebg},frame=single,breaklines=true,columns=fullflexible,keepspaces=true,showstringspaces=false}
\hypersetup{colorlinks=true,linkcolor=dawetblue,urlcolor=dawetblue}
\title{\textcolor{dawetblue}{\textbf{PROJECT DAWET}}\\[.4em]\Large Minggu 5 -- Pra-Assignment 1\\[.3em]\large CSS Layout: Flexbox dan Grid}
\author{Materi Lanjutan HTML \& CSS}
\date{Fase 1 -- Web Foundation}
\begin{document}
\maketitle
\begin{center}\textit{Learn. Build. Understand.}\end{center}
\section{Tujuan Pembelajaran}
Setelah mempelajari materi ini, peserta diharapkan mampu:
\begin{itemize}[nosep]
\item Menjelaskan alasan penggunaan CSS Layout.
\item Memahami container (elemen induk) dan item (elemen anak).
\item Mengatur elemen satu arah menggunakan Flexbox.
\item Menyusun elemen dalam baris dan kolom menggunakan CSS Grid.
\item Memilih teknik layout sesuai kebutuhan halaman web.
\item Menerapkan Flexbox dan Grid pada proyek HTML sederhana.
\end{itemize}
\section{Pengantar CSS Layout}
HTML digunakan untuk membangun struktur halaman, sedangkan CSS digunakan untuk mengatur tampilannya. Ketika halaman semakin kompleks, kita perlu mengatur posisi dan susunan elemen secara lebih terarah.
Contohnya, sebuah header dapat memiliki logo di sebelah kiri, menu navigasi di tengah, dan tombol login di sebelah kanan. CSS Layout membantu menentukan bagaimana elemen-elemen tersebut ditempatkan dan didistribusikan.
\textbf{CSS Layout} adalah pendekatan dalam CSS untuk mengatur bagaimana elemen HTML disusun, ditempatkan, dan menggunakan ruang di dalam halaman.
\textbf{Flexbox} dirancang untuk mengatur elemen dalam satu dimensi, yaitu baris atau kolom. \textbf{CSS Grid} dirancang untuk mengatur elemen dalam dua dimensi, yaitu baris dan kolom sekaligus.
\section{Memahami Container dan Item}
\textbf{Container} adalah elemen induk yang menjadi tempat berkumpulnya elemen lain. \textbf{Item} adalah elemen anak yang berada di dalam container.
\begin{lstlisting}[language=HTML]
<div class="container">
    <div class="item">Satu</div>
    <div class="item">Dua</div>
    <div class="item">Tiga</div>
</div>
\end{lstlisting}
Pada contoh tersebut, container adalah induk, sedangkan tiga elemen item adalah anak. Ketika Flexbox atau Grid diterapkan pada container, item di dalamnya mengikuti aturan layout yang ditentukan.
\section{Flexbox}
Flexbox (Flexible Box Layout) adalah sistem layout CSS untuk menyusun item dalam satu arah: horizontal (baris) atau vertikal (kolom).
\subsection{Mengaktifkan Flexbox}
\begin{lstlisting}[language=CSS]
.container {
    display: flex;
}
\end{lstlisting}
Properti \texttt{display: flex} mengaktifkan Flexbox pada container. Secara default, item disusun dalam satu baris.
\subsection{Mengatur arah item dengan flex-direction}
\begin{lstlisting}[language=CSS]
.container {
    display: flex;
    flex-direction: row;
}
\end{lstlisting}
Nilai umum: \texttt{row} (horizontal), \texttt{column} (vertikal), \texttt{row-reverse} (horizontal terbalik), dan \texttt{column-reverse} (vertikal terbalik).
\subsection{Mengatur posisi item}
\textbf{justify-content} mengatur posisi atau distribusi item pada sumbu utama (\textit{main axis}).
\begin{lstlisting}[language=CSS]
.container {
    display: flex;
    justify-content: center;
}
\end{lstlisting}
Nilai umum: \texttt{flex-start}, \texttt{flex-end}, \texttt{center}, \texttt{space-between}, \texttt{space-around}, dan \texttt{space-evenly}.
\textbf{align-items} mengatur posisi item pada sumbu silang (\textit{cross axis}).
\begin{lstlisting}[language=CSS]
.container {
    display: flex;
    align-items: center;
}
\end{lstlisting}
Sumbu utama mengikuti arah \texttt{flex-direction}. Jika arahnya \texttt{row}, sumbu utama horizontal; jika \texttt{column}, sumbu utama vertikal.
\subsection{Mengatur jarak dengan gap}
\begin{lstlisting}[language=CSS]
.container {
    display: flex;
    gap: 20px;
}
\end{lstlisting}
Properti \texttt{gap} memberikan jarak antar-item secara konsisten tanpa harus menambahkan margin satu per satu.
\section{CSS Grid}
CSS Grid cocok ketika susunan halaman membutuhkan baris dan kolom sekaligus. Contohnya galeri foto, daftar produk, dashboard, serta layout halaman dengan sidebar dan konten utama.
\subsection{Mengaktifkan Grid}
\begin{lstlisting}[language=CSS]
.container {
    display: grid;
}
\end{lstlisting}
\subsection{Membuat kolom}
\begin{lstlisting}[language=CSS]
.container {
    display: grid;
    grid-template-columns: 1fr 1fr 1fr;
}
\end{lstlisting}
Kode tersebut membagi ruang container menjadi tiga kolom sama besar. Penulisan dapat diringkas dengan \texttt{repeat()}:
\begin{lstlisting}[language=CSS]
.container {
    display: grid;
    grid-template-columns: repeat(3, 1fr);
}
\end{lstlisting}
\texttt{1fr} berarti satu bagian dari ruang yang tersedia.
\subsection{Mengatur jarak dan baris}
\begin{lstlisting}[language=CSS]
.container {
    display: grid;
    grid-template-columns: repeat(3, 1fr);
    gap: 20px;
}
\end{lstlisting}
Untuk menentukan tinggi baris, gunakan \texttt{grid-template-rows}:
\begin{lstlisting}[language=CSS]
.container {
    display: grid;
    grid-template-columns: repeat(2, 1fr);
    grid-template-rows: 100px 200px;
}
\end{lstlisting}
Contoh tersebut membuat dua kolom dan dua baris dengan tinggi masing-masing 100px dan 200px.
\section{Flexbox atau Grid?}
\begin{center}
\begin{tabular}{ll}
\toprule
\textbf{Kebutuhan} & \textbf{Pilihan yang sering sesuai}\\
\midrule
Navbar & Flexbox\\
Menyejajarkan tombol & Flexbox\\
Menu horizontal & Flexbox\\
Galeri foto & Grid\\
Daftar produk berbentuk kolom & Grid\\
Dashboard dengan area baris dan kolom & Grid\\
\bottomrule
\end{tabular}
\end{center}
Flexbox dan Grid dapat digunakan bersama. Misalnya, Grid membagi halaman menjadi sidebar dan konten utama, sedangkan Flexbox menyusun tombol di dalam konten.
\section{Praktik Langsung}
Buat dua file dalam satu folder, yaitu \texttt{layout.html} dan \texttt{style.css}.
\subsection{File layout.html}
\begin{lstlisting}[language=HTML]
<!DOCTYPE html>
<html lang="id">
<head>
    <meta charset="UTF-8">
    <meta name="viewport" content="width=device-width, initial-scale=1.0">
    <title>Belajar CSS Layout</title>
    <link rel="stylesheet" href="style.css">
</head>
<body>
    <header class="navbar">
        <h2>DAWET</h2>
        <nav>
            <a href="#">Beranda</a>
            <a href="#">Tentang</a>
            <a href="#">Kontak</a>
        </nav>
    </header>
    <main>
        <h1>Belajar CSS Grid</h1>
        <section class="cards">
            <div class="card">HTML</div>
            <div class="card">CSS</div>
            <div class="card">JavaScript</div>
        </section>
    </main>
</body>
</html>
\end{lstlisting}
\subsection{File style.css}
\begin{lstlisting}[language=CSS]
body {
    font-family: Arial, sans-serif;
    margin: 0;
    padding: 20px;
}
.navbar {
    display: flex;
    justify-content: space-between;
    align-items: center;
}
nav {
    display: flex;
    gap: 20px;
}
.cards {
    display: grid;
    grid-template-columns: repeat(3, 1fr);
    gap: 20px;
}
.card {
    background: #eeeeee;
    padding: 30px;
    text-align: center;
}
\end{lstlisting}
Amati hasilnya: navbar dan menu navigasi menggunakan Flexbox, sedangkan tiga kartu menggunakan Grid. Properti \texttt{gap} mengatur jarak antar-item.
\section{Latihan Mandiri}
\begin{enumerate}
\item Ubah posisi menu navigasi agar berada di sebelah kiri.
\item Buat empat kartu menggunakan Grid.
\item Ubah jarak antarkartu menjadi 30px.
\item Pastikan logo dan menu tetap sejajar secara vertikal.
\item Ubah \texttt{flex-direction} dari \texttt{row} menjadi \texttt{column}, lalu amati perbedaannya.
\end{enumerate}
\section{Evaluasi Pemahaman}
Jawab pertanyaan berikut menggunakan kalimat sendiri:
\begin{enumerate}
\item Apa perbedaan utama Flexbox dan CSS Grid?
\item Apa fungsi container dan item?
\item Apa perbedaan \texttt{justify-content} dan \texttt{align-items}?
\item Mengapa Grid cocok untuk galeri foto atau daftar produk?
\item Berikan satu contoh situasi ketika Flexbox dan Grid digunakan bersama.
\end{enumerate}
\section{Indikator Keberhasilan}
\begin{itemize}
\item Mampu menjelaskan perbedaan Flexbox dan Grid.
\item Mampu menggunakan \texttt{display: flex} dan \texttt{display: grid}.
\item Mampu menyusun navbar dan kartu sederhana.
\item Mampu memodifikasi layout HTML yang sudah dibuat sebelumnya.
\end{itemize}
\section*{Penutup}
Layout bukan sekadar membuat halaman terlihat menarik. Layout membantu mengatur hubungan antar-elemen dan penggunaan ruang di dalam halaman. Pada materi berikutnya, peserta akan mempelajari Responsive Web Design agar halaman tetap nyaman digunakan pada desktop maupun perangkat mobile.
\begin{center}\textit{PROJECT DAWET -- FASE 1\\Learn. Build. Understand.}\end{center}
\end{document}
